% p121_1; source PDF page 121; manually verified editable table.
\begin{tikzpicture}[x=1cm,y=-1cm,inner sep=0pt,outer sep=0pt]
\path[use as bounding box] (0,0) rectangle (16.000,5.120);
\draw[line width=.4pt] (0.000,0.000) rectangle (6.100,1.510);
\node[align=center,text width=5.840cm,font={\linespread{1}\fontsize{10}{13.50}\selectfont\bfseries}] at (3.050,0.755) {案件事实特征};
\draw[line width=.4pt] (6.100,0.000) rectangle (7.200,1.510);
\node[align=center,text width=0.840cm,font={\linespread{1}\fontsize{10}{13.50}\selectfont\bfseries}] at (6.650,0.755) {结论};
\draw[line width=.4pt] (7.200,0.000) rectangle (16.000,0.655);
\node[align=center,text width=8.540cm,font={\linespread{1}\fontsize{10}{13.50}\selectfont\bfseries}] at (11.600,0.328) {名称（术语）};
\draw[line width=.4pt] (7.200,0.655) rectangle (12.100,1.510);
\node[align=center,text width=4.640cm,font={\linespread{1}\fontsize{10}{13.50}\selectfont\bfseries}] at (9.650,1.083) {简洁名称（张明楷教授的观点）};
\draw[line width=.4pt] (12.100,0.655) rectangle (16.000,1.510);
\node[align=center,text width=3.640cm,font={\linespread{1}\fontsize{10}{13.50}\selectfont\bfseries}] at (14.050,1.083) {复杂名称（传统观点）};
\draw[line width=.4pt] (0.000,1.510) rectangle (6.100,3.315);
\node[align=left,text width=5.840cm,font={\linespread{1}\fontsize{10}{13.50}\selectfont}] at (3.050,2.413) {行为对法益没有危险（例如，在荒无人烟的沙漠，甲误以为稻草人是仇人而开枪）};
\draw[line width=.4pt] (6.100,1.510) rectangle (7.200,3.315);
\node[align=center,text width=0.840cm,font={\linespread{1}\fontsize{10}{13.50}\selectfont}] at (6.650,2.413) {无罪};
\draw[line width=.4pt] (7.200,1.510) rectangle (12.100,3.315);
\node[align=center,text width=4.640cm,font={\linespread{1}\fontsize{10}{13.50}\selectfont}] at (9.650,2.413) {不能犯};
\draw[line width=.4pt] (12.100,1.510) rectangle (16.000,3.315);
\node[align=left,text width=3.640cm,font={\linespread{1}\fontsize{10}{13.50}\selectfont}] at (14.050,2.413) {绝对不能犯，不可罚的不能犯};
\draw[line width=.4pt] (0.000,3.315) rectangle (6.100,5.120);
\node[align=left,text width=5.840cm,font={\linespread{1}\fontsize{10}{13.50}\selectfont}] at (3.050,4.218) {行为对法益有危险（例如，甲误将砒霜旁边的白糖当作砒霜，投给乙）};
\draw[line width=.4pt] (6.100,3.315) rectangle (7.200,5.120);
\node[align=center,text width=0.840cm,font={\linespread{1}\fontsize{10}{13.50}\selectfont}] at (6.650,4.218) {有罪};
\draw[line width=.4pt] (7.200,3.315) rectangle (12.100,5.120);
\node[align=center,text width=4.640cm,font={\linespread{1}\fontsize{10}{13.50}\selectfont}] at (9.650,4.218) {未遂犯};
\draw[line width=.4pt] (12.100,3.315) rectangle (16.000,5.120);
\node[align=left,text width=3.640cm,font={\linespread{1}\fontsize{10}{13.50}\selectfont}] at (14.050,4.218) {相对不能犯，可罚的不能犯，不能犯未遂};
\end{tikzpicture}
